\section{From Global State To Local Behaviour}

\begin{proposition}[Nested-state interpretation]
If global conscious states depend on structured dynamical organization, then
local behavioural frames can be interpreted as nested, task-specific regimes
within a wider biological state landscape.
\end{proposition}

\begin{discussion}
This proposition is intentionally conservative. The strong evidence does
\emph{not} say that reading a theorem and falling asleep are the same event.
What it does support is a shared vocabulary: stability, transition,
coordination, and regime change. The book then makes a \textbf{medium}
interpretive move by treating local behavioural frames as smaller-scale regimes
inside that wider landscape. That move is useful because it lets posture,
attention, environment, and cueing be discussed as regime-management variables
rather than as vague mood descriptions.
\end{discussion}

\begin{remark}[What transfers across levels]
The chapter does not import one-to-one mechanisms from neuroscience into
everyday behaviour. It imports a constrained vocabulary only: organized
persistence, transition sensitivity, stability loss, and control-parameter
dependence. That transfer is already enough to improve the manuscript's
behavioural language without pretending to solve the neural realization
problem.
\end{remark}

\begin{discussion}
The bridge can therefore be stated in three steps. First, strong evidence shows
that conscious conditions have organized dynamical structure. Second, this
licenses treating biological persistence as architecture rather than as vague
``mood.'' Third, the manuscript makes a medium-strength interpretive move by
modeling task-level frames as nested regimes that are shaped by that wider
architecture but not identical to it.
\end{discussion}

\begin{remark}[Levels of description]
Wakefulness, attentional coordination, and task engagement are not the same
event. They operate at different descriptive levels. But they are not
independent either. The manuscript works at the behavioural-control layer while
borrowing discipline from the biological layer about what a state, a stable
regime, and a transition are allowed to mean.
\end{remark}

\begin{example}[Same task, different global condition]
Consider the same student facing the same mathematics exercise on two
evenings. After severe sleep loss, the reading frame fragments quickly,
switching cost rises, and distracting cues gain salience. After adequate rest,
the student may still dislike the task, but the regime is less fragile and the
same first step can stabilize. The point is not that biology rigidly dictates
behaviour; it is that global-state organization changes the mechanical
conditions under which local action sequences succeed or fail.
\end{example}
